\documentclass[11pt]{article}
\usepackage[T1]{fontenc}
\usepackage{lmodern,microtype}
\usepackage{amsmath,amssymb,booktabs,enumitem}
\usepackage[letterpaper,margin=1in]{geometry}
\usepackage[colorlinks=true,urlcolor=blue,citecolor=blue]{hyperref}
\newcommand{\F}{\mathbb F}
\newcommand{\Z}{\mathbb Z}
\title{Metaplectic Jacquet--Langlands Transfer, Companion Paper (Rank Two, Wild Place)\\Revision Log: Draft 5 to Draft 6}
\author{Prepared by Claude (Anthropic) for David Joyner}
\date{November 8, 2026}
\begin{document}
\maketitle
\section*{Scope}
Draft 6 responds to ChatGPT's referee report on Draft 5 (five
recommendations).  The main effect is that Theorem \texttt{thm:leveltwo}
no longer depends on any new computation: both of its computational
steps now have hand proofs.  Section 12.5 is rewritten.  Nothing outside
Section 12, the abstract, Appendix B and one line of Section 1 changed.
\section*{Changes}
\begin{enumerate}[leftmargin=*]
\item \textbf{Step 2 by hand} (referee's \S3.1, adopted verbatim in
substance).  $H=Q_3/Q_7$ has $H'=Q_6/Q_7$ of order $3$ and center
$Q_4/Q_7$ of index $9$, so it has $243$ linear and $54$ degree-three
characters, $27$ above each nontrivial $\omega$ of $H'$, and the
isotypic decomposition of the regular character gives
$\sum_{\chi'|_{H'}=3\omega}\chi'=81\,\omega\,\mathbf 1_{H'}$.  Since
$\chi_3$ vanishes off the center of the Heisenberg group, $\chi_3\circ r$
vanishes on the depth-three good classes, and on $H'$ equals
$3\,\omega_3\circ r$ with $D=27$; hence
$(\chi_3\circ r)D=\sum_{\chi'|_{H'}=3\omega_3\circ r}\chi'$, and the
two halves add up to the sum of all $54$.  Script 14 is now
corroboration.
\item \textbf{Step 3 by hand} (new Proposition \texttt{prop:free}).
$\mathrm{Ad}(Q_1)$ acts freely on \emph{all} genuine irreducible
characters of $\hat H=\tilde U_3/\ell(K^7)\ker\tilde\theta_1$.  Proof:
the preimage $\hat A$ of $A=Q_4/Q_7$ is central in $\hat H$ (by
$\omega=0$, Lemma \texttt{lem:H1consequences}(c)), so a genuine
irreducible $\mu$ has a central character $\hat\nu=\lambda_0\alpha$ on
$\hat A$, and $\hat\nu\circ\mathrm{Ad}(t)/\hat\nu=c_t\cdot\omega_\mu([t,\cdot])$
with $c_t=\lambda_0^t/\lambda_0=\zeta^{h_t}$, $h_t\in\Lambda$ nonzero
for $t\notin Q_3$.  For $t\in Q_2$, $[t,A]\subset K^7$ and the ratio is
$c_t|_A$, nontrivial because no nonzero element of $\Lambda$ vanishes on
the $Q_4/Q_6$-coordinates (table in Appendix A).  For $t$ of depth one
with leading coefficient $a$, on $y\in Q_5$ with leading coefficient $b$
the ratio is $\zeta^{2\beta\mathrm{Re}(a\bar b)+2\gamma\mathrm{Im}(a\bar b)}$
($\beta\neq0$ from Lemma \texttt{lem:H1consequences}(b); the commutator
class $2i\,\mathrm{Im}(a\bar b)$ was also checked numerically), a nonzero
functional of $a\bar b$.  Hence no genuine $\mu$ is fixed by any
$t\notin Q_3$.  Script 15 is now corroboration, and writes the
certificate the referee asked for (\texttt{level2\_orbit\_certificate.txt}:
full action tables of the $27$ coset representatives on the $54$
degree-three and $243$ linear genuine characters, explicit stabilizer
test with zero fixed pairs, orbit lists).
\item \textbf{Normalization} (referee's \S3.3).  Section 12 now fixes
$\theta_1:=\mathbf 1$ throughout (allowed by Remark
\texttt{rem:theta-independence}; it is the scripts' normalization), so
$\theta_0=\mathbf 1$ or $\mathrm{sgn}$, and $\lambda_0$, $\hat\Gamma$,
$\hat H$ are all defined in the same normalization; the parenthetical
in Step 1 says how $\ker\tilde\theta_1$ and $\varepsilon_\rho$ are used
when $\theta_\rho=\mathrm{sgn}$.
\item \textbf{Section 12.5 rewritten} (referee's \S4).  (a) The cube
root is a surjection from the good classes of $H_m$ onto $Q_1/Q_{m+1}$,
many-to-one ($219\to27$ for $m=2$), not a bijection.  (b) $\hat H_m$
is defined with $K^{\max(3m+1,6)}$; for $m=1$ it is $\tilde V$, whose
$243$ genuine characters form nine free orbits.  (c) The constituent-wise
``cube-root correspondence'' is withdrawn: the referee's example
($\chi(a+bi)=\zeta^a$, $w=1+i$, $w^5=-w$) shows $\chi\circ r\neq\chi'\circ\pi_3$
for every $\chi'$, and $\Phi_1(\chi\circ r)$ is not a virtual character
for a single additive $\chi\neq1$.  The correct statements are for
$D_8$-invariant characters (restrictions of characters of $D^\times$),
which is all Target B needs; at level two the constituent-wise identity
holds for $\chi_3$ because $\chi_3$ vanishes off the center of the
Heisenberg group.  The section now states the two ingredients
(freeness; $(\Theta_\rho\circ r)D$ a character of $H_m$) that gave the
level theorems, notes that both are open for $m\ge3$, and sketches the
inductive structure of a possible proof of freeness ($\hat A_m$ central
in $\hat H_m$ for $m\ge3$ without hypotheses, since $\ell$ is a
homomorphism on $U_6$).  The term ``cube-root correspondence'' is
removed from the abstract and table.
\item \textbf{Provenance} (referee's \S5).  Appendix B now states that
scripts 13--14 obtain the character table of $Q_3/Q_7$ numerically
(Burnside) and verify it exactly, that orthogonality is a verification
and not a certificate of irreducibility, and that since Draft 6 these
scripts are corroboration only.  The abstract now says ``compatible
with the cover (central character trivial on $\mu_3$)''.  The
introduction's draft range is updated.
\end{enumerate}
\section*{Status}
Theorem \texttt{thm:leveltwo}: [P$^c$], now depending on Target A only
(through $\lambda_0$, $\Lambda$, $e_*$), with no new computational
input.  Proposition \texttt{prop:free}: [P$^c$] for the same reason.
Section 12.5: [O].  Nothing else changed.
\end{document}
